% Chapter 1, typeset from lectures/L01: the README is the opening and the summary, and Appendix A
% and B respectively are the sections, in the same order and under the same subheadings.

\chapter{Combinational logic and first VHDL}
\label{c1:ch}

{\large\itshape From a truth table to a working gate network, in CircuitVerse and in VHDL.\par}

\begin{chapterbox}{In this chapter}
\begin{itemize}
  \item Logic gates and truth tables, and why digital logic tolerates noise that analogue does
    not.
  \item Boolean algebra: reading a sum of products straight out of a truth table.
  \item Building and simulating a gate network by hand in CircuitVerse.
  \item VHDL from the very first session: \code{entity}, \code{architecture}, \code{std_logic} and
    the concurrent assignment.
  \item A complete module, \code{or_gate}, from truth table to synthesisable file.
  \item A brake assistant, derived from a requirement and demonstrated on a DE0-CV board.
\end{itemize}
\end{chapterbox}

This chapter walks the path every circuit in the book walks: from a requirement, by way of a truth
table and a Boolean equation, to a gate network built and simulated by hand, and only then to
VHDL. The order is not a pedagogical detour. It is the working method the rest of the book rests
on, and it is easiest to see on a circuit small enough to take in whole.

Work the chapter in order. \Secref{c1:sec:gates} to \ref{c1:sec:cv} are the circuit;
\secref{c1:sec:vhdl} and \ref{c1:sec:module} are exactly enough VHDL that you can read and write a
gate network; \secref{c1:sec:adas} is the requirement the session builds live, and it is worth
arriving there with an answer of your own to compare against. The exercises in
\secref{c1:sec:exercises} are where you write the code yourself.

\note[Refresher] The gates, the truth tables and the Boolean algebra come with you from digital
electronics. If they need brushing up, they are in sections~1.1--1.4 and 2.1--2.7 of \debook{}, with
exercises and answers, and its appendix~B describes CircuitVerse.

% Sections 1.1-1.6, from lectures/L01/appendix/a_combinational_logic.md (A.1-A.6).

\section{Digital signals and logic gates}
\label{c1:sec:gates}

A digital signal takes one of two values, \code{0} or \code{1}, instead of the continuous range an
analogue signal uses. Anything close enough to \code{0} or \code{1} reads as exactly \code{0} or
\code{1}, so small amounts of electrical noise change nothing. That noise immunity is the reason
nearly all modern computation, from microcontrollers to FPGAs, is digital.

A \textbf{logic gate} has one or more single-bit inputs, a single-bit output and a fixed Boolean
function between them.

\textbf{Combinational logic}, this chapter's subject, produces outputs from the current inputs
alone, with no memory. \textbf{Sequential networks} (\chapref{c3:ch}) depend on the current inputs
\emph{and} on state stored in flip-flops.

Every basic two-input gate, plus the unary \code{NOT}:

\begin{narrowtable}{cc*{6}{c}}
  \thead{A} & \thead{B} & \thead{AND} & \thead{OR} & \thead{NAND} & \thead{NOR} & \thead{XOR} &
    \thead{XNOR} \\
  \midrule
  0 & 0 & 0 & 0 & 1 & 1 & 0 & 1 \\
  0 & 1 & 0 & 1 & 1 & 0 & 1 & 0 \\
  1 & 0 & 0 & 1 & 1 & 0 & 1 & 0 \\
  1 & 1 & 1 & 1 & 0 & 0 & 0 & 1 \\
\end{narrowtable}

\begin{narrowtable}{cc}
  \thead{A} & \thead{NOT A} \\
  \midrule
  0 & 1 \\
  1 & 0 \\
\end{narrowtable}

Worth committing to memory straight from the tables:
\begin{itemize}
  \item \code{NAND}, \code{NOR} and \code{XNOR} are the inverses of \code{AND}, \code{OR} and
    \code{XOR}.
  \item \code{XOR} is \code{1} when the inputs differ. Beyond two inputs it generalises to
    \textbf{odd parity}: \code{1} when an odd number of inputs are \code{1}
    (Exercise~\ref{c1:ex:xor3}).
  \item Every Boolean function can be built from \code{AND}, \code{OR} and \code{NOT}, and also
    from \code{NAND} alone or \code{NOR} alone.
\end{itemize}

% ------------------------------------------------------------------------------------------
\section{Boolean algebra as notation}
\label{c1:sec:boolean}

Three operators, which map straight onto the gates above:

\begin{booktable}{@{}p{5em}L@{}}
  \thead{Gate} & \thead{Notation in Boolean algebra} \\
  \midrule
  AND & multiplication: $A$ \textbf{AND} $B$ is written $AB$ \\
  OR & addition: $A$ \textbf{OR} $B$ is written $A + B$ \\
  NOT & prime: \textbf{NOT} $A$ is written $A'$ \\
\end{booktable}

A function written as an OR of AND terms, for example $X = AB' + CD$, is in \textbf{sum-of-products}
form.

Every truth table can be read off directly as a sum of products:
\begin{enumerate}
  \item Find every row where the output is \code{1}.
  \item For each of them, write an AND term containing every input: plain if the input is \code{1}
    on that row, primed if it is \code{0}.
  \item OR those terms together.
\end{enumerate}

For example:

\begin{narrowtable}{ccc}
  \thead{A} & \thead{B} & \thead{X} \\
  \midrule
  0 & 0 & 0 \\
  0 & 1 & 1 \\
  1 & 0 & 1 \\
  1 & 1 & 0 \\
\end{narrowtable}

The row $AB = 01$ contributes $A'B$ and the row $AB = 10$ contributes $AB'$, which gives
$X = A'B + AB'$. Check it against the table: that is \code{XOR}, so the same network is a single
\code{XOR} gate instead of two \code{AND} gates, two \code{NOT} gates and an \code{OR}.

\subsection*{The laws, as a reminder}
You have met them before. They are here so that the book has a single notation for them and so
that later chapters can refer to them, not because they need teaching:

\begin{booktable}{@{}p{8em}p{9em}L@{}}
  \thead{Law} & \thead{With OR} & \thead{With AND} \\
  \midrule
  Identity & $A + 0 = A$ & $A \cdot 1 = A$ \\
  Dominance & $A + 1 = 1$ & $A \cdot 0 = 0$ \\
  Idempotence & $A + A = A$ & $A \cdot A = A$ \\
  Complement & $A + A' = 1$ & $A \cdot A' = 0$ \\
  Absorption & $A + AB = A$ & $A(A + B) = A$ \\
  De Morgan & $(A + B)' = A'B'$ & $(AB)' = A' + B'$ \\
\end{booktable}

\textbf{De Morgan is the law this book leans on}, less for minimisation than as a rule for how an
inversion moves through a gate: an \code{OR} with both inputs inverted \emph{is} a \code{NAND}, and
an \code{AND} with both inputs inverted \emph{is} a \code{NOR}. That is why a \code{NAND} gate can
build anything, which Exercise~\ref{c1:ex:nand} lets you write in VHDL, and it comes back every
time an active-low signal meets logic written for active-high. In a course where every reset and
every button is active-low, that is often.

The gap between a directly read sum of products and the smallest network is what a \textbf{Karnaugh
map} closes, and \secref{c2:sec:kmap} works through one from start to finish. Not because the
technique is new to you, but because it is the minimised equation that gets written in VHDL, and
\chapref{c8:ch} derives its state machine's next-state logic in exactly that way.

% ------------------------------------------------------------------------------------------
\section{Building and simulating a circuit in CircuitVerse}
\label{c1:sec:cv}

CircuitVerse (\url{https://circuitverse.org/simulator}) is a free, browser-based logic simulator.
Every gate network in this book is built and tested there by hand before it is written in VHDL:
\begin{enumerate}
  \item Open the simulator and start a new project.
  \item Place one \textbf{input} element per input signal and one \textbf{output} element per
    output signal, renamed to match your variables.
  \item Place the gates your equation calls for, from the \emph{Gates} panel.
  \item Wire the network up one gate at a time, following your derived equation. You should
    already know exactly which gates you need before you open the simulator, which is why
    \secref{c1:sec:boolean} comes first.
  \item Toggle every input and check that the output matches your truth table on \textbf{every}
    row.
  \item Save with \code{Project -> Save Online}, or \code{Project -> Download as image/data} for a
    re-importable \file{.cv} file.
\end{enumerate}

\textbf{Step 5 is not optional, and nothing else does it for you.} Three rules make it a check
rather than a formality:
\begin{itemize}
  \item \textbf{Exhaustive, not sampled.} Two inputs are four rows, three are eight. At those
    sizes there is no excuse for sampling. For the sequential networks further on this becomes
    ``step the clock and check at every edge'', the same discipline applied to time.
  \item \textbf{Derive first, simulate second.} Draw the circuit and \emph{then} decide what it
    should do, and you will read your own expectations out of the simulation, which proves
    nothing.
  \item \textbf{A discrepancy is information.} It says the drawing and the equation disagree, and
    that one of them is wrong.
\end{itemize}

Every hardware bug you will chase is found by predicting a behaviour and then checking it. Doing it
by hand on a circuit with four gates is how you learn to do it on one you cannot take in at a
glance.

Once the design is VHDL, a self-checking testbench takes that job over. Every worked example and
every exercise has one handed out with it; you run them rather than write them, which
\secref{c2:sec:testbenches} works through.

% ------------------------------------------------------------------------------------------
\section{Exactly enough VHDL that you can read and write a gate network}
\label{c1:sec:vhdl}

VHDL describes hardware, not a sequence of instructions. Every statement below runs concurrently
and continuously, exactly like the gates it describes. The language arrives as it is needed; this
section is enough for a simple combinational network.

Every VHDL module has two parts:
\begin{itemize}
  \item An \textbf{entity}, the view from outside: the ports, and nothing about what happens
    inside.
  \item An \textbf{architecture}, the implementation behind it.
\end{itemize}

That split is a contract, not ceremony. The entity is everything another design needs in order to
use the module, which is what lets one design instantiate another without reading it, and what lets
you rewrite an implementation without touching what depends on it. The shape is a header beside a
source file, or an interface beside its class; the difference is that VHDL requires it for every
module, with no way to skip the interface or to leak the implementation through it.

Take \code{or_gate}: the inputs \code{a}, \code{b}, the output \code{x}, and $x = a + b$.

\begin{vhdlcode}
entity or_gate is
    port(a, b: in std_logic;
         x   : out std_logic);
end entity;
\end{vhdlcode}

\lecfig[0.62]{lectures/L01/appendix/images/or_gate_entity.png}{The entity \code{or_gate}: the view
from outside, and nothing about what happens inside.}{c1:fig:entity}

\code{std_logic} is the signal type for a single bit. It is not built into VHDL but comes from the
package \code{std_logic_1164}, which is why every file in this book opens with:

\begin{vhdlcode}
library ieee;
use ieee.std_logic_1164.all;
\end{vhdlcode}

It has nine values instead of two, because real hardware needs more than true and false:
\begin{itemize}
  \item \code{'0'} and \code{'1'}: driven low and driven high. The only two you will write in this
    book.
  \item \code{'Z'}: high impedance, nothing is driving the wire at all.
  \item \code{'U'}: uninitialised, what a signal holds before anything has assigned it a value.
  \item \code{'X'}: unknown, what a simulator shows when two things drive the same signal and
    disagree.
  \item \code{'-'}, \code{'W'}, \code{'L'}, \code{'H'}: don't-care values and weakly driving
    values, rarely written by hand.
\end{itemize}

Recognising the seven you will not write still matters: they are how a simulator tells you
something is wrong. A signal left sitting at \code{'U'} or \code{'X'} is a bug, not a value.

The architecture fills the box in:

\begin{vhdlcode}
architecture behaviour of or_gate is
begin
    x <= a or b;
end architecture;
\end{vhdlcode}

\lecfig[0.62]{lectures/L01/appendix/images/or_gate_arch.png}{The architecture \code{or_gate}: the
implementation behind the entity.}{c1:fig:arch}

Together they make up the complete module:

\lecfig[0.62]{lectures/L01/appendix/images/or_gate_module.png}{The module \code{or_gate}: entity and
architecture together.}{c1:fig:module}

\code{x <= a or b;} is \secref{c1:sec:boolean}'s Boolean algebra almost verbatim: VHDL's
\code{and}, \code{or}, \code{not} and \code{xor} work directly on \code{std_logic}, so every
equation you derived by hand becomes VHDL with hardly any translation at all.

The assignment is \textbf{concurrent}. It does not run once; it holds continuously, like a wire
between real gates.

The keyword \code{signal} declares an internal wire inside an architecture, as opposed to a port
visible from outside. You use one in Exercise~\ref{c1:ex:adas} to name an intermediate result
instead of nesting it in place, and \chapref{c2:ch} leans on the same idea throughout.

% ------------------------------------------------------------------------------------------
\section{The complete module: \code{or_gate}}
\label{c1:sec:module}

\Secref{c1:sec:vhdl}'s parts together make a complete, synthesisable file. This is the whole of
\file{or_gate.vhd}, and it is a real design that synthesises and runs on the DE0-CV:

\vhdlfile{lectures/L01/or_gate/or_gate.vhd}

\begin{itemize}
  \item The \code{library}/\code{use} clauses pull in \code{std_logic}.
  \item \code{entity} declares the outside: two inputs, one output.
  \item \code{architecture} implements the inside, as a single concurrent assignment.
\end{itemize}

That is the smallest complete VHDL design there is, and the shape never changes. Every module in
this book, all the way up to the CAN controller in \chapref{c17:ch}, is the same
\code{library}/\code{entity}/\code{architecture} skeleton. What grows is the architecture body,
never the frame.

\paragraph{Checking your work}
\code{or_gate} comes with an \file{or_gate_tb.vhd}, which drives all four input combinations and
checks the output:

\begin{shell}
cd lectures/L01/or_gate
ghdl -a --std=93 or_gate.vhd or_gate_tb.vhd
ghdl -e --std=93 or_gate_tb
ghdl -r --std=93 or_gate_tb --assert-level=error --stop-time=10ms
\end{shell}

No output beyond a closing note means every check passed. Testbenches are treated properly in
\secref{c2:sec:testbenches}; until then, take these three commands as the way you confirm that a
module works.

% ------------------------------------------------------------------------------------------
\section{A brake assistant, from requirement to circuit}
\label{c1:sec:adas}

The session builds a different circuit from \code{or_gate}, live, so that you get to see the same
path walked twice: once here in writing, and once on something you have not already read the answer
to.

The requirement: \textbf{the car brakes if the driver brakes, or if the assistance system decides
to.} The assistance system decides to when either detector sees something, unless it is at the same
time reporting a fault, in which case it is not trusted to brake at all.

\begin{booktable}{@{}p{8em}p{4em}L@{}}
  \thead{Port} & \thead{Direction} & \thead{Meaning} \\
  \midrule
  \code{driver_brake} & in & The driver is on the brake pedal. \\
  \code{sensor} & in & The distance sensor sees an obstacle. \\
  \code{radar} & in & The radar sees an obstacle. \\
  \code{error} & in & The assistance system reports a fault. \\
  \code{engine_brake} & out & Apply the brakes. \\
\end{booktable}

Three sentences together determine the whole structure:
\begin{itemize}
  \item either detector is enough on its own, and neither of them outweighs the other.
  \item a fault raises no alarm; it \emph{prevents} the assistance system from being able to brake.
  \item the driver can always brake, even when the assistance system is broken and doing nothing.
\end{itemize}

The last is a safety requirement, and it is the one to hold on to as you derive the network: the
pedal is not an input to the assistance logic, it goes past it. Work out what changes if the pedal
is instead fed through the fault logic, and you have found the bug the structure exists to prevent.

\subsection*{The truth table}
Sixteen rows, in the order a four-bit counter would produce them. \code{adas_brake} is not a port;
it is the assistance system's own decision, shown as an intermediate column so that you can see
where the structure comes from.

\begin{narrowtable}{cccccc}
  \thead{\code{driver_brake}} & \thead{\code{sensor}} & \thead{\code{radar}} &
    \thead{\code{error}} & \thead{\code{adas_brake}} & \thead{\code{engine_brake}} \\
  \midrule
  0 & 0 & 0 & 0 & 0 & 0 \\
  0 & 0 & 0 & 1 & 0 & 0 \\
  0 & 0 & 1 & 0 & 1 & 1 \\
  0 & 0 & 1 & 1 & 0 & 0 \\
  0 & 1 & 0 & 0 & 1 & 1 \\
  0 & 1 & 0 & 1 & 0 & 0 \\
  0 & 1 & 1 & 0 & 1 & 1 \\
  0 & 1 & 1 & 1 & 0 & 0 \\
  1 & 0 & 0 & 0 & 0 & 1 \\
  1 & 0 & 0 & 1 & 0 & 1 \\
  1 & 0 & 1 & 0 & 1 & 1 \\
  1 & 0 & 1 & 1 & 0 & 1 \\
  1 & 1 & 0 & 0 & 1 & 1 \\
  1 & 1 & 0 & 1 & 0 & 1 \\
  1 & 1 & 1 & 0 & 1 & 1 \\
  1 & 1 & 1 & 1 & 0 & 1 \\
\end{narrowtable}

The lower half is \code{1} throughout: as soon as \code{driver_brake} is set, nothing else changes
the answer. That is what ``overrides everything'' looks like in a table.

\subsection*{Four answers to bring to the session}
Everything above is the specification. The derivation, the gate network and the VHDL code are what
gets built live, and it is worth arriving there with an answer of your own to compare against
rather than one you have already read. So, from the table above:
\begin{itemize}
  \item Read a sum-of-products equation for \code{engine_brake} straight out of it, the same way
    \secref{c1:sec:boolean} did.
  \item Simplify it by hand. The requirement is three sentences long, so the minimal network is
    considerably smaller than what the row-by-row reading gives you. Work out how much smaller
    before you are told.
  \item Count the gates your simplified equation needs, and write the number down.
  \item Sketch the network, and decide which single gate the safety argument above rests on.
\end{itemize}

Bring those four answers to the session. Being wrong about one of them is worth more than not
having guessed, because you then know exactly which step to look at again when the live derivation
reaches it. It is \secref{c1:sec:cv}'s ``derive first, simulate second'' applied to a circuit you
have not built yet, and it is the habit the rest of the book is built on. Afterwards you write the
VHDL code yourself, in Exercise~\ref{c1:ex:adas}, and a testbench decides whether your equations and
the truth table above agree.

\subsection*{Why there is a testbench for it}
Sixteen combinations are few enough to check by hand once. They are not few enough to check every
time someone edits the file, which is the situation every module in this book is in from here on.

A \textbf{testbench} is a second VHDL file that does the checking. It is not part of the design and
never reaches the FPGA: it drives the inputs through every case, works out what the output should
have been, compares and reports. For \code{adas}, which is written alongside it during the session,
that is a loop over the sixteen combinations restating the braking rule independently.

\emph{Independently} is the important word. A testbench that read its expected answer out of the
design would agree with it by construction, even when the design is wrong. Restating the
requirement separately is what makes it a check.

That is why nearly every directory in the course repository has a \file{_tb.vhd} beside the module:
it is how you can be handed a module, change it, and know within a second whether you have broken
it. \Secref{c2:sec:testbenches} works through how they are run; you are not asked to write one
anywhere in this book.

% Section 1.7, from lectures/L01/README.md: the goals, the self-check questions and the look
% ahead.

\section{Summary}
\label{c1:sec:review}

\needspace{6\baselineskip}
\subsection*{What you should be able to do now}
\begin{itemize}
  \item Derive a sum of products from a truth table, and then build and simulate the network in
    CircuitVerse.
  \item Apply De Morgan to move an inversion through a gate, and use that to realise a function in
    a gate set that lacks the operator you started from.
  \item Read and write a combinational VHDL module, an \code{entity} with \code{std_logic} ports
    and an \code{architecture} that drives its output, and say why VHDL keeps the two apart.
  \item Explain why a concurrent assignment describes a wire rather than a statement that runs
    once, and why a signal left sitting at \code{'U'} or \code{'X'} is a fault and not a value.
  \item Run a testbench handed out with an exercise and decide whether your module passed it.
\end{itemize}

\needspace{6\baselineskip}
\subsection*{Questions to test yourself with}
\begin{enumerate}
  \item What is the difference between an \code{entity} and an \code{architecture}, and why does
    VHDL separate them?
  \item Why is \code{x <= a or b;} not ``a statement that runs once''? What does it describe
    instead?
  \item In the brake assistant the driver's pedal is ORed in last, instead of going through the
    fault-handling logic. What breaks if you wire it the other way round?
\end{enumerate}

\needspace{6\baselineskip}
\subsection*{Looking ahead}
\Chapref{c2:ch} goes on with:
\begin{itemize}
  \item Karnaugh maps: finding the \emph{smallest} gate network, not just a correct one.
  \item \code{std_logic_vector}, \code{process} and \code{case}.
  \item Submodules, demonstrated on a two-digit hexadecimal display driven by eight switches.
  \item Self-checking testbenches: verifying your own VHDL without a board.
\end{itemize}

% Section 1.8, from lectures/L01/appendix/b_exercises.md.

\section{Exercises}
\label{c1:sec:exercises}

\begin{aside}
\textbf{How to check your work.} Exercise~\ref{c1:ex:sop} is a short pen-and-paper warm-up on the
notation the book uses throughout; it needs neither VHDL, GHDL nor a board, and it states what a
good answer covers.

From Exercise~\ref{c1:ex:gates} onwards, every exercise that asks you to write a VHDL module comes
with a self-checking testbench under \file{lectures/L01/exercises/}. Write your module in its
directory \file{exercises/<module>/}, with the entity name and the \textbf{port order} the exercise
states, and then run it with GHDL \textemdash{} see \secref{c2:sec:testbenches} for the three
commands.

No FPGA board is needed for any exercise. The Quartus synthesis and board programming steps are
demonstrated during the session; your job afterwards is to get the VHDL right, and the testbench is
how you confirm it.
\end{aside}

% ------------------------------------------------------------------------------------------
\exerciseset{Truth tables and gates}

\exercise{A function that depends on less than it looks like it does}{Circuit}
\label{c1:ex:sop}
A function of two inputs \code{A}, \code{B} and one output \code{X} has the truth table below:

\begin{narrowtable}{ccc}
  \thead{A} & \thead{B} & \thead{X} \\
  \midrule
  0 & 0 & 1 \\
  0 & 1 & 0 \\
  1 & 0 & 1 \\
  1 & 1 & 0 \\
\end{narrowtable}

\xpart{a} Derive the sum-of-products equation for \code{X} straight out of the truth table
(\secref{c1:sec:boolean}). Two rows have $X = 1$, so you get two AND terms.

\xpart{b} Simplify it algebraically. Factor out what the two terms share, and apply $A + A' = 1$
and $1 \cdot B' = B'$ from \secref{c1:sec:boolean}. You should end up with a single literal.

\xpart{c} Your answer to \textbf{b)} says something about \code{A}. What? Which single gate from
the table in \secref{c1:sec:gates} computes the function, and what happens to the input \code{A}
when you draw it?

\xpart{d} Build that single-gate version in CircuitVerse and confirm that it reproduces every row
of the truth table above. Leave the input \code{A} unconnected, or simply leave it out of the
drawing; \textbf{b)} is the justification for doing so.

% ------------------------------------------------------------------------------------------
\exerciseset{Entities and architectures}

\exercise{\code{and_gate} and \code{nand_gate}}{VHDL}
\label{c1:ex:gates}
Write two small gates as complete VHDL modules, \code{and_gate} and \code{nand_gate}.

Neither gate is the point, and you should not expect either of them to be hard. The point is that
this is where an entity and an architecture stop being something you read about in
\secref{c1:sec:vhdl} and become something you write, and that the two halves are worth writing in
that order.

Both entities have the same ports:

\begin{narrowtable}{lll}
  \thead{Port} & \thead{Direction} & \thead{Type} \\
  \midrule
  \code{a} & in & \code{std_logic} \\
  \code{b} & in & \code{std_logic} \\
  \code{x} & out & \code{std_logic} \\
\end{narrowtable}

\xpart{a} Write the \code{entity} declaration for \code{and_gate}, and nothing else. In VHDL a
module's interface is something you can write down and reason about before any implementation
exists at all, which is the habit this part exists to build. Identify which ports are required, the
direction of each, and the type of each.

\xpart{b} Add an \code{architecture} that implements the AND gate with a single concurrent signal
assignment.

\xpart{c} Now write \code{nand_gate}: the same interface with an inverted output, using VHDL's
\code{nand} operator directly. Do \textbf{not} write \code{not (a and b)}.

The two are equivalent as Boolean algebra, and on an FPGA they make no difference at all, since
both end up in the same lookup table. It matters on the kind of hardware Exercise~\ref{c1:ex:nand}
below is about, where NAND is the gate you physically have and everything else must be built from
it. Writing what you mean, and letting the tools decide what it costs, is the habit; reaching past
an operator the language already gives you is what to avoid.

\modulefig[0.52]{lectures/L01/appendix/images/and_gate.png}
\modulefig[0.52]{lectures/L01/appendix/images/nand_gate.png}

\begin{selfcheck}
\textbf{Self-check.} Name your entities \code{and_gate} and \code{nand_gate}, each with the ports
\code{a}, \code{b} (in) and \code{x} (out), declared in that order; their testbenches are in
\file{exercises/and_gate/} and \file{exercises/nand_gate/}, and each checks all four rows of the
corresponding truth table.
\end{selfcheck}

\exercise{\code{xor3}: odd parity with three inputs}{VHDL}
\label{c1:ex:xor3}
An XOR gate has three inputs, \code{a}, \code{b}, \code{c}, and one output, \code{x}. The output is
\code{1} when an odd number of the inputs are \code{1} (three-input XOR, or the parity function).

\xpart{a} Write the truth table for \code{x} over all eight combinations of \code{a}, \code{b},
\code{c}, and confirm that it matches the description ``an odd number of \code{1}s''.

\xpart{b} Implement the gate in VHDL, as a module of its own with an entity whose ports are three
\code{std_logic} inputs and one \code{std_logic} output, and an architecture that drives \code{x}
with a single concurrent assignment. \Secref{c1:sec:vhdl} shows the entity/architecture pattern and
notes that VHDL's \code{xor} operator works directly on \code{std_logic}; \secref{c1:sec:module}
shows a complete module from start to finish.

\xpart{c} Build the same gate in CircuitVerse and confirm that it reproduces every row of your
truth table.

\note[Hint]\code{xor} chains from left to right, so the whole architecture body is a single line,
\code{x <= a xor b xor c;}, with no need for either an internal \code{signal} or intermediate
gates.

\modulefig[0.52]{lectures/L01/appendix/images/xor3.png}

\begin{selfcheck}
\textbf{Self-check.} Name your VHDL entity \code{xor3}, with the ports \code{a}, \code{b},
\code{c} (in) and \code{x} (out), declared in that order; its testbench is in
\file{exercises/xor3/}, and \secref{c1:sec:module} lists the three commands that run it.
\end{selfcheck}

\exercise{\code{majority3}: a triple-redundant vote}{VHDL}
\label{c1:ex:majority3}
A \textbf{majority gate} has three inputs, \code{a}, \code{b}, \code{c}, and one output, \code{x}.
The output is \code{1} when \emph{at least two} of the three inputs are \code{1}.

This is the circuit behind a triple-redundant vote: three sensors report the same measurement, and
the system acts on the answer two of them agree on, so that a single broken sensor cannot decide
anything on its own.

\xpart{a} Write the truth table for \code{x} over all eight combinations of \code{a}, \code{b},
\code{c}.

\xpart{b} Derive the sum-of-products equation for \code{x} straight out of the truth table
(\secref{c1:sec:boolean}):
\begin{itemize}
  \item Four rows have $x = 1$, so you get four AND terms.
  \item Every term contains all three inputs, primed where the input is \code{0} on its row.
\end{itemize}

\xpart{c} Implement the gate in VHDL, as a module of its own:
\begin{itemize}
  \item An entity with three \code{std_logic} inputs and one \code{std_logic} output.
  \item An architecture that drives \code{x} with a single concurrent assignment, translating your
    four AND terms directly (\secref{c1:sec:vhdl}).
\end{itemize}

\xpart{d} Build the same network in CircuitVerse and confirm that it reproduces every row of your
truth table.

\note[Hint]Write the equation out on paper before you open either the editor or CircuitVerse. Four
three-input AND terms plus a four-input OR gate is a lot of wiring to correct after the fact.

\note[Looking ahead]Count the gates your equation calls for, and then convince yourself by
inspection that $x = ab + ac + bc$ computes the same function with three two-input AND gates
instead of four three-input ones. Sum of products gave you a \emph{correct} network, not the
\emph{smallest}; finding the smaller one systematically, rather than by staring at it, is what
\chapref{c2:ch} is for.

\modulefig[0.52]{lectures/L01/appendix/images/majority3.png}

\begin{selfcheck}
\textbf{Self-check.} Name your VHDL entity \code{majority3}, with the ports \code{a}, \code{b},
\code{c} (in) and \code{x} (out), declared in that order; its testbench is in
\file{exercises/majority3/} and checks all eight rows. Either equation passes \textemdash{} the
testbench checks the function, not which network you chose.
\end{selfcheck}

\exercise{\code{half_adder}: the first module with two outputs}{VHDL}
\label{c1:ex:half_adder}
Write a \textbf{half adder}: \code{sum} and \code{carry} from two single bits. You know the
circuit; the reason it is here is that it is the book's first module with \textbf{two outputs}, and
that is a VHDL question rather than a logic question.

\xpart{a} Implement it, with two concurrent assignments and no internal \code{signal}.

Two outputs do not mean two entities, two architectures or a process. Every output port is driven
by its own concurrent assignment, and the two are not steps that happen in an order: they describe
two separate pieces of wire, both live all the time. If you catch yourself reaching for a process
in order to ``return'' two values, that instinct is the software instinct, and
\secref{c1:sec:vhdl}'s ``a concurrent assignment describes a wire'' is the correction.

\xpart{b} Build it in CircuitVerse and confirm both outputs, so that you have seen the two-output
shape as a drawing before you meet it as an entity with several \code{out} ports in
\chapref{c2:ch}.

\note[Looking ahead]This is the cell the counter in \chapref{c6:ch} is built from. The four-bit
adder there is four of these chained carry-to-carry, and the counter is that adder with its output
fed back to its input, which is why the wraparound needs no logic at all.

\modulefig[0.56]{lectures/L01/appendix/images/half_adder.png}

\begin{selfcheck}
\textbf{Self-check.} Name your entity \code{half_adder}, with the ports \code{a}, \code{b} (in) and
\code{sum}, \code{carry} (out), declared in that order; its testbench is in
\file{exercises/half_adder/} and checks \code{sum} and \code{carry} separately, so that a design
that has one right and the other wrong is told which is which.
\end{selfcheck}

% ------------------------------------------------------------------------------------------
\exerciseset{Universal gates}

\exercise{\code{or_from_nand}}{VHDL}
\label{c1:ex:nand}
\Secref{c1:sec:gates} states that every Boolean function can be built from \code{NAND} gates
alone. This exercise has you prove it to yourself for the three operators you have been using.

\xpart{a} Derive the \code{OR} case on paper before you build anything, since De Morgan
(\secref{c1:sec:boolean}) gives it to you directly. Start from $(A'B')'$ and apply the law; you
should arrive at $A + B$. Now read that expression back as a circuit: what is $A'B'$ with an
inversion on its output, and how many \code{NAND} gates is the whole thing?

\xpart{b} Build each of the following in CircuitVerse from \code{NAND} gates alone, checking each
against its truth table in \secref{c1:sec:gates} as you go:
\begin{itemize}
  \item \code{NOT a}: one NAND gate is enough, if you feed the same signal to both its inputs.
  \item \code{a AND b}: one NAND gate, followed by the \code{NOT} you just built.
  \item \code{a OR b}: the network you derived in \textbf{a)}. Check your gate count against what
    you predicted.
\end{itemize}

\xpart{c} Write the \code{OR} version in VHDL, as a module of its own:
\begin{itemize}
  \item An entity with two \code{std_logic} inputs and one \code{std_logic} output.
  \item An architecture that drives \code{x} with a single concurrent assignment.
  \item Use VHDL's \code{nand} operator alone: no \code{or}, no \code{and}, no \code{not}
    anywhere in the expression.
\end{itemize}

\xpart{d} Explain why this matters in practice. Bear in mind that a chip manufacturing process has
to physically implement every gate type it offers, and that \code{NAND} is the cheapest gate to
build in CMOS.

\note[Hint]VHDL's \code{nand} operator does not chain the way \code{and} and \code{or} do:
\code{a nand b nand c} is not valid. Parenthesise every NAND explicitly, which is what you want
here anyway, since every pair of parentheses is a physical gate in your CircuitVerse drawing.

\modulefig[0.52]{lectures/L01/appendix/images/or_from_nand.png}

\begin{selfcheck}
\textbf{Self-check.} Name your VHDL entity \code{or_from_nand}, with the ports \code{a}, \code{b}
(in) and \code{x} (out), declared in that order; its testbench is in
\file{exercises/or_from_nand/} and checks it against all four rows of the \code{OR} truth table.
\end{selfcheck}

% ------------------------------------------------------------------------------------------
\exerciseset{Putting it all together}

\exercise{The brake assistant in VHDL}{VHDL}
\label{c1:ex:adas}
Write the brake assistant from \secref{c1:sec:adas} in VHDL.

\Secref{c1:sec:adas} had you derive its equations, count its gates and sketch its network
\emph{before} the session, and the session then built it live. This is where you write it yourself,
and where a testbench decides whether your reading of the requirement matches the one in the truth
table.

The entity must have:

\begin{booktable}{@{}p{8em}p{3.4em}p{5.6em}L@{}}
  \thead{Port} & \thead{Dir.} & \thead{Type} & \thead{Description} \\
  \midrule
  \code{driver_brake} & in & \code{std_logic} & The driver is on the brake pedal. \\
  \code{sensor} & in & \code{std_logic} & The distance sensor sees an obstacle. \\
  \code{radar} & in & \code{std_logic} & The radar sees an obstacle. \\
  \code{error} & in & \code{std_logic} & The assistance system reports a fault. \\
  \code{engine_brake} & out & \code{std_logic} & Apply the brakes. \\
\end{booktable}

\xpart{a} Implement it from \textbf{your own} equations rather than from what was on the screen. If
the two disagree, that disagreement is the most useful thing this exercise can give you, and the
testbench says which of them the truth table sides with.

Use two concurrent assignments and one internal \code{signal}:
\begin{itemize}
  \item The signal holds the assistance system's own decision, before the driver's pedal is
    weighed in.
  \item This is the book's first exercise with an internal signal, and it is the reason the keyword
    exists: a name for an intermediate result, so that every line still reads as the sentence in
    the requirement it came from.
\end{itemize}

\xpart{b} Now break it deliberately, in the specific way \secref{c1:sec:adas} warns about. Feed
\code{driver_brake} \emph{through} the fault logic instead of letting it go past:

\begin{vhdlcode}
engine_brake <= (driver_brake or sensor or radar) and (not error);
\end{vhdlcode}

It is one gate simpler, and it looks reasonable. Run the testbench against it and answer:
\begin{itemize}
  \item Which input combination does it fail on first, and what happens physically in the car at
    that moment?
  \item How many of the sixteen rows does it get wrong, and what do they all have in common?
  \item The rows that fail are exactly the ones the safety requirement was written for. What does
    that tell you about the value of a truth table derived from a requirement, as opposed to one
    read off a circuit you have already built?
\end{itemize}

\modulefig[0.72]{lectures/L01/appendix/images/adas.png}

\begin{selfcheck}
\textbf{Self-check.} Name your entity \code{adas}, with the inputs \code{driver_brake},
\code{sensor}, \code{radar}, \code{error} and the output \code{engine_brake}, declared in that
order; its testbench is in \file{exercises/adas/} and checks all sixteen rows of
\secref{c1:sec:adas}'s truth table.
\end{selfcheck}

